\documentclass[12pt, a4paper]{ctexart}
\usepackage{fontspec}
\usepackage{xcolor}
\usepackage{hyperref}
\usepackage{geometry}
\usepackage{enumitem}
\usepackage{parskip}
\usepackage{titlesec}
\usepackage{tcolorbox}
\tcbuselibrary{breakable}
\tcbset{autoparskip}

\usepackage{setspace}
\usepackage{listings}
\usepackage{amssymb}

\geometry{left=2.5cm, right=2.5cm, top=2.5cm, bottom=2.5cm}

\setmainfont{Times New Roman}
\setCJKmainfont{SimSun}

\hypersetup{
    colorlinks=true,
    linkcolor=blue!70!black,
    urlcolor=cyan,
    pdftitle={JDBC-逐题精讲解义},
    pdfauthor={专升本-fifine老师}
}

\definecolor{myblue}{RGB}{30,100,200}
\definecolor{lightblue}{RGB}{230,240,255}
\definecolor{darkblue}{RGB}{0,50,120}

\newtcolorbox{bluebox}[1][]{
    breakable,
    colback=lightblue,
    colframe=myblue,
    coltitle=darkblue,
    fonttitle=\bfseries,
    title={#1},
    boxrule=0.8pt,
    arc=4pt,
    left=6pt,
    right=6pt,
    top=6pt,
    bottom=6pt,
    before skip=8pt,
    after skip=8pt
}

\usepackage{etoolbox}
\BeforeBeginEnvironment{bluebox}{\par\vfill}%
\AfterEndEnvironment{bluebox}{\par\vfill}%

\titleformat{\section}
  {\normalfont\Large\bfseries\color{myblue}}{\thesection}{1em}{}
\titlespacing*{\section}{0pt}{3.5ex plus 1ex minus .2ex}{1.5ex plus .2ex}

\title{\textcolor{myblue}{\textbf{JDBC - 逐题精讲解义}}}
\author{专升本-fifine老师 教学组}
\date{2026年03月05日}

\lstset{
    language=Java,
    basicstyle=\ttfamily\small,
    keywordstyle=\color{blue},
    commentstyle=\color{green!60!black},
    stringstyle=\color{orange},
    numbers=left,
    numberstyle=\tiny\color{gray},
    frame=single,
    breaklines=true,
    columns=fullflexible,
    showstringspaces=false
}

\newcommand{\code}[1]{\texttt{#1}}

\begin{document}

\maketitle
\thispagestyle{empty}
\newpage

\section{JDBC}

\begin{enumerate}[label=\textbf{\arabic*.}]

	\item \textbf{以下关于Java中JDBC的说法错误的是}
	      \begin{enumerate}[label=\Alph*.]
		      \item 用于连接数据库
		      \item 不同的数据库有不同的驱动
		      \item 执行SQL语句通过Statement对象
		      \item 不需要关闭连接资源
	      \end{enumerate}

	      \begin{bluebox}{答案与深度解析}
		      \textbf{答案：D. 不需要关闭连接资源}

		      \textbf{【核心认知框架】}
		      \begin{itemize}[label={}, leftmargin=*, nosep]
			      \item \textbf{题目本质}：考查JDBC的基本使用规范和资源管理
			      \item \textbf{关键区分点}：JDBC资源占用问题，必须显式关闭
			      \item \textbf{命题人意图}：测试对JDBC资源管理的理解
		      \end{itemize}

		      \textbf{【选项原理深度拆解】}

		      \begin{itemize}[leftmargin=*, nosep, label={}]
			      \item[A] \textbf{用于连接数据库 — 正确}
			            \begin{itemize}[label={}, leftmargin=*, nosep]
				            \item \textbf{功能说明}：JDBC提供Java连接数据库的标准API
				            \item \textbf{连接对象}：Connection表示与数据库的连接
				            \item \textbf{代码示例}：
				                  \begin{lstlisting}[language=Java]
import java.sql.*;

public class JDBCDemo {
    public static void main(String[] args) {
        // 数据库连接信息
        String url = "jdbc:mysql://localhost:3306/test";
        String user = "root";
        String password = "password";

        try {
            // 加载驱动（JDBC 4.0+可省略）
            Class.forName("com.mysql.cj.jdbc.Driver");

            // 建立连接
            Connection conn = DriverManager.getConnection(url, user, password);
            System.out.println("数据库连接成功！");

            // 执行SQL操作...

            conn.close();
        } catch (ClassNotFoundException e) {
            System.out.println("驱动未找到");
        } catch (SQLException e) {
            e.printStackTrace();
        }
    }
}
// JDBC用于建立和管理数据库连接
				      \end{lstlisting}
				            \item \textbf{连接参数}：
				                  \begin{itemize}[label={}, nosep]
					                  \item URL：数据库地址（jdbc:mysql://host:port/db）
					                  \item User：数据库用户名
					                  \item Password：数据库密码
				                  \end{itemize}
			            \end{itemize}

			      \item[B] \textbf{不同的数据库有不同的驱动 — 正确}
			            \begin{itemize}[label={}, leftmargin=*, nosep]
				            \item \textbf{驱动概念}：每个数据库厂商提供对应的JDBC驱动实现
				            \item \textbf{常见驱动}：
				                  \begin{itemize}[label={}, nosep]
					                  \item MySQL：com.mysql.cj.jdbc.Driver
					                  \item Oracle：oracle.jdbc.driver.OracleDriver
					                  \item PostgreSQL：org.postgresql.Driver
					                  \item SQL Server：com.microsoft.sqlserver.jdbc.SQLServerDriver
					                  \item SQLite：org.sqlite.JDBC
				                  \end{itemize}
				            \item \textbf{代码示例}：
				                  \begin{lstlisting}[language=Java]
// MySQL驱动
String mysqlDriver = "com.mysql.cj.jdbc.Driver";
String mysqlURL = "jdbc:mysql://localhost:3306/test";

// Oracle驱动
String oracleDriver = "oracle.jdbc.driver.OracleDriver";
String oracleURL = "jdbc:oracle:thin:@localhost:1521:orcl";

// PostgreSQL驱动
String pgDriver = "org.postgresql.Driver";
String pgURL = "jdbc:postgresql://localhost:5432/test";

// 加载不同的驱动连接不同的数据库
Class.forName(mysqlDriver);
Connection mysqlConn = DriverManager.getConnection(mysqlURL);

Class.forName(oracleDriver);
Connection oracleConn = DriverManager.getConnection(oracleURL);
// 不同数据库使用不同的驱动实现
				      \end{lstlisting}
				            \item \textbf{JDBC 4.0+特性}：自动加载驱动，无需显式Class.forName()
			            \end{itemize}

			      \item[C] \textbf{执行SQL语句通过Statement对象 — 正确}
			            \begin{itemize}[label={}, leftmargin=*, nosep]
				            \item \textbf{Statement体系}：
				                  \begin{itemize}[label={}, nosep]
					                  \item Statement：执行静态SQL
					                  \item PreparedStatement：执行参数化SQL（推荐）
					                  \item CallableStatement：执行存储过程
				                  \end{itemize}
				            \item \textbf{代码示例}：
				                  \begin{lstlisting}[language=Java]
Connection conn = DriverManager.getConnection(url, user, password);

// 使用Statement
Statement stmt = conn.createStatement();
String sql1 = "SELECT * FROM users";
ResultSet rs = stmt.executeQuery(sql1);

String sql2 = "UPDATE users SET age = age + 1";
int rows = stmt.executeUpdate(sql2);

// 使用PreparedStatement（推荐）
String sql3 = "SELECT * FROM users WHERE id = ?";
PreparedStatement pstmt = conn.prepareStatement(sql3);
pstmt.setInt(1, 100);  // 设置参数（从1开始）
ResultSet rs2 = pstmt.executeQuery();

// 执行存储过程
String sql4 = "{call get_user_info(?, ?)}";
CallableStatement cstmt = conn.prepareCall(sql4);
cstmt.setInt(1, 100);
cstmt.registerOutParameter(2, Types.VARCHAR);
cstmt.execute();
String name = cstmt.getString(2);
// JDBC通过Statement/PreparedStatement等执行SQL
				      \end{lstlisting}
				            \item \textbf{方法区别}：
				                  \begin{itemize}[label={}, nosep]
					                  \item executeQuery()：执行SELECT，返回ResultSet
					                  \item executeUpdate()：执行INSERT/UPDATE/DELETE，返回影响行数
					                  \item execute()：执行任意SQL，返回boolean表示是否有ResultSet
				                  \end{itemize}
			            \end{itemize}

			      \item[D] \textbf{不需要关闭连接资源 — 错误（本题答案）}
			            \begin{itemize}[label={}, leftmargin=*, nosep]
				            \item \textbf{错误原因}：JDBC资源占用数据库连接，必须显式关闭
				            \item \textbf{需要关闭的资源}：
				                  \begin{itemize}[label={}, nosep]
					                  \item ResultSet：结果集
					                  \item Statement：语句对象
					                  \item PreparedStatement：预处理语句
					                  \item Connection：数据库连接
				                  \end{itemize}
				            \item \textbf{资源关闭顺序}：ResultSet → Statement → Connection
				            \item \textbf{代码反证}：
				                  \begin{lstlisting}[language=Java]
import java.sql.*;

public class NotClosingDemo {
    public static void main(String[] args) {
        try {
            Connection conn = DriverManager.getConnection(url, user, password);
            Statement stmt = conn.createStatement();
            ResultSet rs = stmt.executeQuery("SELECT * FROM users");

            // 使用结果集...

            // 未关闭资源！
            // conn.close();
            // stmt.close();
            // rs.close();
        } catch (SQLException e) {
            e.printStackTrace();
        }
        // 问题：连接资源一直占用，导致连接池耗尽
        // 数据库连接是有限资源，不关闭会导致连接泄漏
    }
}
// 不关闭连接资源会导致连接泄漏
					      \end{lstlisting}
				            \item \textbf{正确做法（传统方式）}：
				                  \begin{lstlisting}[language=Java]
Connection conn = null;
Statement stmt = null;
ResultSet rs = null;
try {
    conn = DriverManager.getConnection(url, user, password);
    stmt = conn.createStatement();
    rs = stmt.executeQuery("SELECT * FROM users");

    while (rs.next()) {
        // 处理数据
    }
} catch (SQLException e) {
    e.printStackTrace();
} finally {
    // 按顺序关闭资源
    if (rs != null) {
        try { rs.close(); } catch (SQLException e) {}
    }
    if (stmt != null) {
        try { stmt.close(); } catch (SQLException e) {}
    }
    if (conn != null) {
        try { conn.close(); } catch (SQLException e) {}
    }
}
// 使用finally块确保资源关闭
					      \end{lstlisting}
				            \item \textbf{推荐做法（try-with-resources）}：
				                  \begin{lstlisting}[language=Java]
// Java 7+ try-with-resources自动关闭资源
try (Connection conn = DriverManager.getConnection(url, user, password);
     Statement stmt = conn.createStatement();
     ResultSet rs = stmt.executeQuery("SELECT * FROM users")) {

    while (rs.next()) {
        // 处理数据
    }
} catch (SQLException e) {
    e.printStackTrace();
}
// 自动关闭Connection、Statement、ResultSet
// 只需要实现AutoCloseable接口的资源
					      \end{lstlisting}
				            \item \textbf{连接池场景}：使用连接池时，归还连接而不是关闭
				                  \begin{lstlisting}[language=Java]
import javax.sql.DataSource;

DataSource dataSource = ...; // 获取数据源
Connection conn = null;
try {
    conn = dataSource.getConnection();  // 从池获取
    Statement stmt = conn.createStatement();
    // 执行SQL...
    stmt.close();
} catch (SQLException e) {
    e.printStackTrace();
} finally {
    if (conn != null) {
        conn.close();  // 归还池，不是真正关闭
    }
}
// 连接池环境下，close()将连接归还池中复用
					      \end{lstlisting}
			            \end{itemize}
		      \end{itemize}

		      \textbf{【应背诵知识点】}
		      \begin{itemize}[label={}, nosep]
			      \item JDBC用于连接和操作数据库
			      \item 不同数据库使用不同驱动（MySQL/Oracle等）
			      \item 通过Statement/PreparedStatement执行SQL
			      \item 必须关闭资源（ResultSet→Statement→Connection）
			      \item 推荐使用try-with-resources自动关闭
			      \item 口诀：JDBC连接数据库，不同驱动各异样，Statement执行SQL，资源关闭不能忘
		      \end{itemize}

		      \textbf{【命题趋势与实战策略】}
		      \textbf{考场秒杀}：看到"不需要关闭连接"→错，必须关闭资源
		      \textbf{高频陷阱}：忘记关闭ResultSet、Statement、Connection
	      \end{bluebox}

\end{enumerate}

\end{document}
